\section{Pointers and References}

\subsection{Section Overview}
Pointers are one of the most powerful and feared features of C++. They allow for direct memory manipulation, which is essential for low-level programming, dynamic memory allocation, and building efficient data structures. In this section, we'll demystify pointers, learn how to use them safely, and also cover references, which are a simpler and safer alternative in many cases.

\subsection{What is a Pointer?}
A pointer is a variable that stores the memory address of another variable. Instead of holding a value (like an \texttt{int} or a \texttt{double}), a pointer ``points to'' the location where a value is stored.

\subsection{Declaring Pointers}
You declare a pointer by specifying the type of the data it will point to, followed by an asterisk \texttt{*}, and then the pointer's name.
\begin{lstlisting}
	int *int_ptr;       // points to an integer
	double *double_ptr; // points to a double
	std::string *str_ptr; // points to a string
\end{lstlisting}
\begin{remark}
An uninitialized pointer points to a random memory address. This is dangerous. Always initialize your pointers. A common practice is to initialize pointers to \texttt{nullptr} if they don't point to a valid address.
\end{remark}

\subsection{Storing Addresses in Pointers}
The address-of operator (\texttt{\&}) gives you the memory address of a variable. You can use this to make a pointer point to a variable.
\begin{lstlisting}
	int num = 10;
	int *p_num = &num; // p_num now holds the address of num

	std::cout << "Value of num: " << num << std::endl;
	std::cout << "Address of num: " << &num << std::endl;
	std::cout << "Value of p_num: " << p_num << std::endl; // will be same as &num
\end{lstlisting}

\subsection{Dereferencing a Pointer}
Dereferencing a pointer means accessing the value at the memory address the pointer is holding. The dereference operator is also an asterisk \texttt{*}.
\begin{lstlisting}
	int num = 10;
	int *p_num = &num;

	std::cout << "Value pointed to by p_num: " << *p_num << std::endl; // prints 10

	*p_num = 20; // changes the value of num to 20
	std::cout << "New value of num: " << num << std::endl; // prints 20
\end{lstlisting}

\subsection{Dynamic Memory Allocation}
Pointers are crucial for dynamic memory allocation, which is allocating memory on the heap at runtime. The \texttt{new} keyword allocates memory and returns a pointer to it. The \texttt{delete} keyword frees the memory.
\begin{lstlisting}
	int *p_int = nullptr;
	p_int = new int; // allocates memory for one int
	*p_int = 100;

	std::cout << "Value: " << *p_int << std::endl; // 100

	delete p_int; // frees the allocated memory

	// Dynamically allocating an array
	int *p_array = new int[10]; // allocates memory for 10 ints
	p_array[0] = 5;

	delete[] p_array; // use delete[] for arrays
\end{lstlisting}
\begin{remark}
For every \texttt{new}, there must be a corresponding \texttt{delete}. For every \texttt{new[]}, there must be a corresponding \texttt{delete[]}. Failing to free memory leads to memory leaks. In modern C++, smart pointers are used to manage dynamic memory automatically.
\end{remark}

\subsection{Arrays and Pointers}
The name of an array is actually a pointer to its first element.
\begin{lstlisting}
	int scores[] {100, 95, 88};
	int *p_scores = scores; // p_scores points to the first element

	std::cout << *scores << std::endl; // 100
	std::cout << *(scores + 1) << std::endl; // 95
	std::cout << *(p_scores + 2) << std::endl; // 88
\end{lstlisting}

\subsection{Pointer Arithmetic}
You can perform arithmetic operations on pointers. Incrementing a pointer moves it to the next element of its type, not just the next byte.
\begin{lstlisting}
	int scores[] {100, 95, 88};
	int *p_scores = scores;

	std::cout << *p_scores << std::endl; // 100
	p_scores++;
	std::cout << *p_scores << std::endl; // 95
\end{lstlisting}

\subsection{const and Pointers}
\texttt{const} can be used with pointers in three ways:
\begin{lstlisting}
	// 1. Pointer to a constant: The data pointed to cannot be changed.
	const int *p1;
	int const *p2; // same as above

	// 2. Constant pointer: The pointer itself cannot be changed to point to something else.
	int *const p3 = &some_int;

	// 3. Constant pointer to a constant: Both are constant.
	const int *const p4 = &some_int;
\end{lstlisting}

\subsection{Passing Pointers to Functions}
Passing a pointer to a function allows the function to access and modify the original data. This is more efficient than passing large objects by value.
\begin{lstlisting}
	void double_it(int *p_num) {
		*p_num *= 2;
	}
	int main() {
		int x = 10;
		double_it(&x); // pass the address of x
		std::cout << x << std::endl; // 20
	}
\end{lstlisting}

\subsection{Returning a Pointer from a Function}
You can return a pointer from a function, but you must be careful not to return a pointer to a local variable, as it will be destroyed when the function exits. It's usually for returning dynamically allocated memory.
\begin{lstlisting}
	int* create_array(size_t size, int init_value = 0) {
		int *new_array = new int[size];
		for (size_t i=0; i<size; ++i)
			new_array[i] = init_value;
		return new_array;
	}
\end{lstlisting}

\subsection{Potential Pointer Pitfalls}
\begin{itemize}
	\item \textbf{Dangling pointers:} Pointers that point to memory that has been freed.
	\item \textbf{Null pointer dereference:} Attempting to dereference a \texttt{nullptr}.
	\item \textbf{Memory leaks:} Forgetting to \texttt{delete} memory allocated with \texttt{new}.
\end{itemize}

\subsection{What is a Reference?}
A reference is an alias for an already existing variable. Once a reference is initialized to a variable, it cannot be changed to refer to another variable. References are generally safer and easier to use than pointers.
\begin{lstlisting}
	int num = 10;
	int &ref_num = num; // ref_num is an alias for num

	ref_num = 20; // num is now 20
	std::cout << num << std::endl; // 20
\end{lstlisting}
We have already used references when passing arguments to functions (pass-by-reference).

\subsection{Coding Exercise 30: Find the maximum element in an array using pointers}
Write a function that takes a pointer to the beginning of an integer array and its size, and returns a pointer to the maximum element.
\begin{lstlisting}
	#include <iostream>

	int* find_max(int *arr, int size) {
		if (size == 0) return nullptr;
		int *max_ptr = arr;
		for (int i = 1; i < size; ++i) {
			if (*(arr + i) > *max_ptr) {
				max_ptr = arr + i;
			}
		}
		return max_ptr;
	}

	int main() {
		int arr[] = {10, 5, 25, 15, 30};
		int *max_el = find_max(arr, 5);
		std::cout << "Max element is: " << *max_el << std::endl;
		return 0;
	}
\end{lstlisting}

\subsection{Coding Exercise 31: Reverse an array using pointers}
Write a function that reverses an array in place using two pointers.
\begin{lstlisting}
	#include <iostream>

	void reverse_array(int *arr, int size) {
		int *start = arr;
		int *end = arr + size - 1;
		while (start < end) {
			int temp = *start;
			*start = *end;
			*end = temp;
			start++;
			end--;
		}
	}

	int main() {
		int arr[] = {1, 2, 3, 4, 5};
		reverse_array(arr, 5);
		for (int i = 0; i < 5; ++i) {
			std::cout << arr[i] << " ";
		}
		std::cout << std::endl;
		return 0;
	}
\end{lstlisting}

\subsection{Coding Exercise 32: Reverse a string using pointers}
Write a function to reverse a C-style string in place using pointers.
\begin{lstlisting}
	#include <iostream>
	#include <cstring>

	void reverse_string(char *str) {
		char *start = str;
		char *end = str + strlen(str) - 1;
		while (start < end) {
			char temp = *start;
			*start = *end;
			*end = temp;
			start++;
			end--;
		}
	}

	int main() {
		char my_str[] = "hello";
		reverse_string(my_str);
		std::cout << my_str << std::endl;
		return 0;
	}
\end{lstlisting}

\subsection{L-values and R-values}
An \textbf{l-value} (locator value) represents an object that has an identifiable location in memory (i.e., an address). You can take the address of an l-value. Variables are l-values. An \textbf{r-value} is a temporary value that does not have a persistent memory address. Literals (like \texttt{10}) are r-values. Pointers can only point to l-values.

\subsection{Section Challenge}
Write a function that takes a pointer to an array of integers and its size, and returns a pointer to a new dynamically allocated array that contains the elements of the original array in reverse order. The original array should not be modified.
\subsection{Section Challenge --- Solution}
\begin{lstlisting}
	#include <iostream>

	int* reverse_copy(const int *arr, int size) {
		if (size == 0) return nullptr;
		int *reversed_arr = new int[size];
		for (int i = 0; i < size; ++i) {
			reversed_arr[i] = arr[size - 1 - i];
		}
		return reversed_arr;
	}

	int main() {
		int arr[] = {1, 2, 3, 4, 5};
		int *new_arr = reverse_copy(arr, 5);

		std::cout << "Original array: ";
		for(int i=0; i<5; ++i) std::cout << arr[i] << " ";
		std::cout << std::endl;

		std::cout << "Reversed copy: ";
		for(int i=0; i<5; ++i) std::cout << new_arr[i] << " ";
		std::cout << std::endl;

		delete[] new_arr; // Don't forget to free the memory!

		return 0;
	}
\end{lstlisting}

\subsection{Quiz 9: Section 12 Quiz}
\begin{enumerate}
	\item What does a pointer variable store?
	\begin{enumerate}
		\item[a)] An integer value.
		\item[b)] A floating-point value.
		\item[c)] A memory address.
		\item[d)] A character.
	\end{enumerate}

	\item What is the operator used to get the address of a variable?
	\begin{enumerate}
		\item[a)] \texttt{*}
		\item[b)] \texttt{\&}
		\item[c)] \texttt{->}
		\item[d)] \texttt{.}
	\end{enumerate}

	\item What is dereferencing a pointer?
	\begin{enumerate}
		\item[a)] Deleting the pointer.
		\item[b)] Getting the address stored in the pointer.
		\item[c)] Accessing the value at the address stored in the pointer.
		\item[d)] Changing the pointer to \texttt{nullptr}.
	\end{enumerate}
\end{enumerate}

\textbf{Answers:} 1. (c), 2. (b), 3. (c)
